\documentclass[12pt, oneside]{book}
\usepackage{../style/mypackages}
\usepackage{../style/abbr}





\begin{document}
\frontmatter
\title{{\Huge{\textbf{Differential Manifolds}}}}
\maketitle

\dominitoc % 初始化minitoc
\pagenumbering{Roman}
\tableofcontents % 生成内容目录 


\mainmatter


\chapter{Manifolds}

\section{Manifolds} % Manifolds

\subsection{Topological Manifolds}  % Topological Manifolds
\begin{definition}
    Suppose that $M$ is a topological space. We say that $M$ is a \textbf{topological manifold} of dimension $n$ or a topological n-manifold if it has the following properties:
    \begin{enumerate}[label=(\roman*)]
        \item $M$ is a Hausdorff space
        \item $M$ is second-countable
        \item $M$ is locally Euclidean of dimension $n$: each point $p$ of $M$ has a neighborhood $U$ that is homeomorphic to an open subset $\widehat{U}$ of $\mathbb{R}^n$. $\left(U,\varphi\right)$ called \textbf{coordinate chart}.
    \end{enumerate}
    Given a chart $(U, \varphi)$, we call the set $U$ a \textbf{coordinate domain}, and $\varphi$ is called a \textbf{local coordinate map}.
\end{definition}





\subsection{Topological Properties of Manifolds}

\begin{definition}
    Let $M$ be a topological space.
    \begin{enumerate}
        \item
              A collection $\mathcal{X}$ of subsets of $M$ is said to be \textbf{locally finite} if each point of $M$ has a neighborhood that intersects at most finitely many of the sets in $\mathcal{X}$.

        \item
              Given a cover $\mathcal{U}$ of $M$, another cover $\mathcal{V}$ is called a \textbf{refinement} of $\mathcal{U}$ if for each $V \in \mathcal{V}$ there exists some $U \in \mathcal{U}$ such that $V \subseteq U$.

        \item
              We say that $M$ is \textbf{paracompact} if every open cover of $M$ admits an open, locally finite refinement.
    \end{enumerate}
\end{definition}

\begin{proposition}
    Suppose $M$ is a topological manifold, then
    \begin{enumerate}
        \item
              local compact

        \item
              $\sigma$-compact

        \item
              $M$ is Lindelof's : any open cover of $M$ has a countable subcover.
        \item
              paracompact
    \end{enumerate}
\end{proposition}





\begin{theorem}[Paracompact]
    Every topological manifold is paracompact. In fact, given a topological manifold $M$, an open cover $\mathcal{X}$ of $M$, and any basis $\mathfrak{B}$ for the topology of $M$, there exists a countable, locally finite open refinement of $\mathcal{X}$ consisting of elements of $\mathcal{B}$.

    Proof:
    Given $M, \mathcal{X}$, and $\mathcal{B}$ as in the hypothesis of the theorem, let $\left(K_j\right)_{j=1}^{\infty}$ be an exhaustion of $M$ by compact sets (Proposition A.60). For each $j$, let $V_j=K_{j+1} \backslash$ $\operatorname{Int} K_j$ and $W_j=\operatorname{Int} K_{j+2} \backslash K_{j-1}$ (where we interpret $K_j$ as $\varnothing$ if $j<1$ ). Then $V_j$ is a compact set contained in the open subset $W_j$. For each $x \in V_j$, there is some $X_x \in \mathcal{X}$ containing $x$, and because $\mathcal{B}$ is a basis, there exists $B_x \in \mathcal{B}$ such that $x \in B_x \subseteq X_x \cap W_j$. The collection of all such sets $B_x$ as $x$ ranges over $V_j$ is an open cover of $V_j$, and thus has a finite subcover. The union of all such finite subcovers as $j$ ranges over the positive integers is a countable open cover of $M$ that refines $\mathcal{X}$. Because the finite subcover of $V_j$ consists of sets contained in $W_j$, and $W_j \cap W_{j^{\prime}}=\varnothing$ except when $j-2 \leq j^{\prime} \leq j+2$, the resulting cover is locally finite.
\end{theorem}











\subsection{Smooth Structures and Smooth Manifolds}
\begin{definition}
    Let $M$ be a topological $n$-manifold.
    \begin{enumerate}
        \item
              Two charts $(U, \varphi)$ and $(V, \psi)$ are said to be \textbf{smoothly compatible} if either $U \cap V=\varnothing$ or the two \textbf{transition map}
              \begin{equation*}
                  \psi \circ \varphi^{-1}: \varphi(U \cap V) \rightarrow \psi(U \cap V),\quad
                  \varphi \circ \psi^{-1}: \psi(U \cap V) \rightarrow \varphi(U \cap V)
              \end{equation*}
              are $C^\infty$.

        \item
              We define an \textbf{smooth atlas} $\mathcal{A}$ for $M$ to be a collection of charts $\left\{\left(U_\alpha, \varphi_\alpha\right)\right\}$ whose domains cover $M$ and any two charts in $\mathcal{A}$ are smoothly compatible with each other.

        \item
              A smooth atlas $\mathcal{A}$ on $M$ is \textbf{maximal} if it is not properly contained in any larger smooth atlas.
              If $M$ is a topological manifold, a \textbf{smooth structure} on $M$ is a maximal smooth atlas.

              A $C^\infty$ manifold is a pair $(M, \mathcal{A})$, where $M$ is a topological manifold and $\mathcal{A}$ is a \textbf{smooth structure} on $M$.
    \end{enumerate}
\end{definition}

Although the $C^\infty$ compatibility of charts is clearly reflexive and symmetric, it is
not transitive. However, we have the following lemma:

\begin{lemma}
    Let $\left\{\left(U_\alpha, \phi_\alpha\right)\right\}$ be an atlas on a locally Euclidean space. If two charts $(V, \psi)$ and $(W, \sigma)$ are both compatible with the atlas $\left\{\left(U_\alpha, \phi_\alpha\right)\right\}$, then they are compatible with each other.
\end{lemma}

\begin{proposition}
    Let $M$ be a topological manifold.
    \begin{enumerate}
        \item
              Every smooth atlas $\mathcal{A}$ for $M$ is contained in a unique maximal smooth atlas $\overline{\mathcal{A}}$, called the \textbf{smooth structure determined by $\mathcal{A}$}. Indeed, $\overline{\mathcal{A}}$ denote the set of all charts that
              are smoothly compatible with every chart in $\mathcal{A}$.

        \item
              Two smooth atlases for $M$ determine the same smooth structure if and only if
              their union is a smooth atlas.
    \end{enumerate}
\end{proposition}




\subsection{Manifold with Boundary}
Denote by $\mathbb{R}^n_+$ or $\mathbb{H}^n$ the subset of $\mathbb{R}^n$ defined by
\begin{equation*}
    \mathbb{R}^n_+=\left\{\left(x^1, \ldots, x^n\right) \in \mathbb{R}^n: x^n \geq 0\right\}
\end{equation*}



\begin{definition}
    An \textbf{n-dimensional topological manifold with boundary} is a second-countable Hausdorff space $M$ in which every point $p$ has a neighborhood $U$ homeomorphic $\varphi$ to an open subset of $\mathbb{R}^n_+$.
    \begin{enumerate}
        \item
              The open subset $U \subseteq M$ together with a map $\varphi: U \rightarrow \mathbb{R}^n_+$ that is a homeomorphism onto an open subset of $\mathbb{R}^n_+$  will be called a \textbf{chart} for $M$.
        \item
              We will call $(U, \varphi)$ an \textbf{interior chart} if $\varphi(U)$ is an open subset of $\mathbb{R}^n$, and a \textbf{boundary chart} if  $\varphi(U) \cap \partial \mathbb{R}^n_+ \neq \varnothing$.
        \item
              A point $p \in M$ is called \textbf{interior point} if it is in the domain of some interior chart.

              A point $p \in M$ is called \textbf{boundary point} if it is in the domain of a boundary chart that sends $p$ to $\partial \mathbb{R}^n_+$.
    \end{enumerate}
    The boundary of $M$ (the set of all its boundary points) is denoted by $\partial M$; similarly, its interior, the set of all its interior points, is denoted by $\operatorname{Int} M$.
\end{definition}


\begin{theorem}
    Each point of $M$ is either an interior point or a boundary point, but not both.
\end{theorem}

\begin{corollary}
    Let $M$ be a topological $n$-manifold with boundary.
    \begin{enumerate}
        \item
              $\Int M$ is an open subset of $M$ and a topological n-manifold without boundary.

        \item
              $\partial M$ is a closed subset of $M$ and a topological $(n-1)$-manifold without boundary.

        \item
              $M$ is a topological manifold if and only if $\partial M=\varnothing$.

        \item
              If $n=0$, then $\partial M=\varnothing$ and $M$ is a $0$ -manifold.
    \end{enumerate}
\end{corollary}






\subsubsection{Smooth Manifold with Boundary}
Recall that a map from an arbitrary subset $A \subseteq \mathbb{R}^n$ to $\mathbb{R}^k$ is said to be smooth if in a neighborhood of each point of $A$ it admits an extension to a smooth map defined on an open subset of $\mathbb{R}^n$.
\begin{definition}
    Let $M$ be a topological manifold with boundary. A \textbf{smooth structure} for $M$ is defined to be a maximal smooth atlas $\mathcal{A}$, a collection of charts whose domains cover $M$ and whose transition maps (and their inverses) are smooth.
    With such a structure, $\left(M,\mathcal{A}\right)$ is called a \textbf{smooth manifold with boundary}.
\end{definition}



\section{Smooth Maps on a Manifold}


\begin{definition}
    Let $M$ be a manifold, $k$ is a nonnegative integer, and $f: M \rightarrow \mathbb{R}^k$ is any function. We say that $f$ is \textbf{smooth at $p$} if there exists a chart $(U, \varphi)$ for $M$ whose domain contains $p$ and such that the composite function
    \begin{equation*}
        f \circ \varphi^{-1} : \varphi(U)\rightarrow
        \mathbb{R}^k
    \end{equation*}
    is smooth at $\varphi(p)$.
    The function $f \circ \varphi^{-1}$ is called the \textbf{local coordinate representation} of $f$.
    The function $f$ is said to be $C^{\infty}$ on $M$ if it is $C^{\infty}$ at every point of $M$.
    \begin{remark}
        The definition of the smoothness of a function $f$ at a point is independent of the chart $(U,\varphi)$.
    \end{remark}
\end{definition}

\begin{proposition}
    Let $M$ be a manifold of dimension $n$, and $f: M \rightarrow \mathbb{R}$ a real-valued function on $M$. The following are equivalent:
    \begin{enumerate}
        \item
              The function $f: M \rightarrow \mathbb{R}$ is $C^{\infty}$.

        \item
              The manifold $M$ has an atlas such that for every chart $\left(U, \phi\right)$ in the atlas, $f \circ \phi^{-1}: \mathbb{R}^n \supset \phi(U) \rightarrow \mathbb{R}$ is $C^{\infty}$.

        \item
              For every chart $(V, \psi)$ on $M$, the function $f \circ \psi^{-1}: \mathbb{R}^n \supset \psi(V) \rightarrow \mathbb{R}$ is $C^{\infty}$.
    \end{enumerate}
\end{proposition}
\begin{proposition}[Smoothness in terms of components]
    Let $M$ be a manifold.
    A vector-valued function $F: M \rightarrow \mathbb{R}^m$ is $C^{\infty}$ if and only if its component functions $F^1, \ldots, F^m: M \rightarrow \mathbb{R}$ are all $C^{\infty}$.
\end{proposition}




\begin{definition}
    Let $M, N$ be smooth manifolds, and let $f: M \rightarrow N$ be a continuous map. We say that $f$ is a \textbf{smooth map} if for every $p \in M$, there exist smooth charts $(U, \varphi)$ containing $p$ and $(V, \psi)$ containing $f(p)$ (assume that $f(U) \subseteq V$ without generality) and the composite map
    \begin{equation*}
        \tilde{f}=\psi \circ f \circ \varphi^{-1} :\varphi(f^{-1}(V)\cap U)\rightarrow \psi(V)
    \end{equation*}
    is smooth. We call $\tilde{f}=\psi \circ f \circ \varphi^{-1}$ the \textbf{coordinate representation
        of $f$ with respect to the given coordinates}.
\end{definition}

\begin{proposition}
    Let $N$ and $M$ be smooth manifolds, and $F: N \rightarrow M$ a continuous map. The following are equivalent:
    \begin{enumerate}
        \item
              The map $F: N \rightarrow M$ is $C^{\infty}$.

        \item
              There are atlases $\mathfrak{U}$ for $N$ and $\mathfrak{V}$ for $M$ such that for every chart $\left(U, \phi\right)$ in $\mathfrak{U}$ and $(V, \psi)$ in $\mathfrak{V}$, the map
              \begin{equation*}
                  \psi \circ F \circ \phi^{-1}: \phi\left(U \cap F^{-1}(V)\right) \rightarrow \mathbb{R}^m
              \end{equation*}
              is $C^{\infty}$.

        \item
              For every chart $\left(U, \phi\right)$ on $N$ and $\left(V, \psi\right)$ on $M$, the map
              \begin{equation*}
                  \psi \circ F \circ \phi^{-1}: \phi\left(U \cap F^{-1}(V)\right) \rightarrow \mathbb{R}^m
              \end{equation*}
              is $C^{\infty}$.
    \end{enumerate}
\end{proposition}




\begin{proposition}[Composition of $C^\infty$ maps]
    Let $M, N$, and $P$ be smooth manifolds with or without boundary.
    \begin{enumerate}
        \item
              Every constant map $M \rightarrow N$ is smooth.
        \item  The identity map of $M$ is smooth.
        \item
              If $U \subseteq M$ is an open submanifold with or without boundary, then the inclusion map $U \hookrightarrow M$ is smooth.
        \item
              If $F: M \rightarrow N$ and $G: N \rightarrow P$ are smooth, then so is $G \circ F: M \rightarrow P$.
    \end{enumerate}
\end{proposition}


\begin{definition}
    If $M$ and $N$ are smooth manifolds with or without boundary, a \textbf{diffeomorphism} from $M$ to $N$ is a smooth bijective map $F: M \rightarrow N$ that has a smooth inverse. We say that $M$ and $N$ are \textbf{diffeomorphic} if there exists a diffeomorphism between them. Sometimes this is symbolized by $M \approx N$.
\end{definition}
\begin{proposition}
    Let $U$ be an open subset of a manifold $M$ of dimension $n$ and $f: U \rightarrow f(U) \subset \mathbb{R}^n$ be a maps.
    Then $f$ is a diffeomorphism onto its image if and only if $(U, f)$ is a coordinate chart in the differentiable structure of $M$.
\end{proposition}



\section{Partition of unity}

\subsection{Topological Preliminaries}

\begin{definition}
    Let $X$ be a topological space, if there exists $X_j$ such that
    \begin{enumerate}[label=(\roman*)]
        \item For each $j$, the closure $\overline{X}_j$ is compact.

        \item For each $j$, $ \overline{X}_j \subset X_{j+1}$.

        \item  $M=\bigcup_j X_j$.
    \end{enumerate}
    The subsets $ \left\{X_j\right\}$ described is called an \textbf{exhaustion} of $M$.
\end{definition}



\begin{definition}
    A real-valued continuous function $f$ on $X$ is called an \textbf{exhausion function} for $X$ if for any $c \in \mathbb{R}$, the sublevel set $f^{-1}\left((\infty, c]\right)$ is compact.
\end{definition}

\begin{lemma}[Lindelof's]
    Let $X$ be a second countable space, then every open cover $\mathcal{A}$ has a countable subcover.
\end{lemma}

\begin{theorem}
    If $X$ is a second countable, locally compact and Hausdorff space (thus a manifold), then there exists a exhaustion of $X$.
\end{theorem}

\begin{definition}
    Let $\mathcal{A}$ be a open cover of the space $X$, then a open cover $\mathcal{B}$ of $X$ is said to be a \textbf{refinement} of $\mathcal{A}$ (or is said to refine $\mathcal{A}$ ) if for each element $B$ of $\mathcal{B}$, there is an element $A$ of $\mathcal{A}$ containing $B$.
\end{definition}
\begin{definition}
    A space $X$ is \textbf{paracompact} if every open cover $\mathcal{A}$ of $X$ has a locally finite open      refinement $\mathcal{B}$ of $\mathcal{A}$.
\end{definition}

\begin{theorem}
    Let $M$ be any topological manifold. For any open cover $\mathcal{U}=\left\{U_\alpha\right\}$ of $M$, one can find two countable family of open covers $\mathcal{V}=\left\{V_j\right\}$ and $\mathcal{W}=\left\{W_j\right\}$ of $M$ so that
    \begin{enumerate}[label=(\roman*)]
        \item $\mathcal{W}$ is a locally finite open refinement of $\mathcal{U}$.
        \item For each $j$, $ \overline{V}_j$ is compact and $\overline{V}_j \subset W_j$.

    \end{enumerate}

    Proof.
    Let $\left\{X_j\right\}$ be a exhaustion of $M$.  For each $p \in M$, there is an $j$ and an $\alpha(p)$ so that $p \in \overline{X}_{j+1} \backslash X_j$ and $p \in U_{\alpha(p)}$.
    Since $M$ is locally Euclidean, one can always choose open neighborhoods $V_p, W_p$ of $p$ so that $\overline{V}_p$ is compact and
    \begin{equation*}
        p
        \in
        V_p
        \subset
        \overline{V}_p
        \subset
        W_p
        \subset
        U_{\alpha(p)} \cap\left(X_{j+2} \backslash \overline{X}_{j-1}\right)
    \end{equation*}
    Now for each $j$, since the "stripe" $\overline{X}_{j+1} \backslash X_j$ is compact, one can choose finitely many points $p_1^j, \cdots, p_{k_j}^j$ so that $V_{p_1^j}, \cdots, V_{p_{k_j}^j}$ is an open cover of $\overline{X}_{j+1} \backslash X_j$.
    Denote all these $V_{p_k^j}$ 's by $V_1, V_2, \cdots$, and the corresponding $W_{p_k^j}$ 's by $W_1, W_2, \cdots$. Then $\mathcal{V}=\left\{V_k\right\}$ and $\mathcal{W}=\left\{W_k\right\}$ are open covers of $M$ that satisfies all the conditions.
\end{theorem}


\begin{corollary}
    Manifolds is paracompact, $\sigma$-compact
\end{corollary}




\subsection{Existence}
\begin{lemma}
    There exists $f_1,f_2, f_3$ such that
\end{lemma}

\begin{theorem}[Bump function on manifold]
    Let $M$ be a smooth manifold, $K \subset M$ is a compact subset, and $U \subset M$ an open subset that contains $K$.
    Then there is a $\varphi \in C^{\infty}(M)$ so that $K \prec f \prec U$ ($0 \leq \varphi \leq 1, \varphi \equiv 1$ on $K$ and $\operatorname{supp}(\varphi) \subset U$).

    Proof. For each $q \in K$, there is a chart $\left(\varphi_q, U_q, V_q\right)$ near $q$ so that $U_q \subset U$ and $V_q$ contains the open ball $B_3(0)$ in $\mathbb{R}^n$.
    Let $U'_q=\varphi_q^{-1}\left(B_1(0)\right)$, and let
    \begin{equation*}
        f_q(p)= \begin{cases}f_3\left(\varphi_q(p)\right) & , p \in U_q    \\
             0                            & , p \notin U_q\end{cases}
    \end{equation*}
    Then $f_q \in C^{\infty}(M), \operatorname{supp}\left(f_q\right) \subset U_q \subset U$ and $f_1 \equiv 1$ on ${U}'_q$.

    Now the family of open sets $\left\{U'_q\right\}$ is an open cover of $K$.
    Since $K$ is compact, there is a finite sub-cover $\left\{{U}'_{q_1}, \cdots, U'_{q_n}\right\}$. Let \begin{equation*}
        \psi=\sum_{i=1}^n f_{q_i}
    \end{equation*}
    Then $\psi$ is a smooth and compactly supported function on $M$ so that $\psi \geq 1$ on $K$ and $\operatorname{supp}(\psi) \subset U$.
    It follows that the function $\varphi(p)=f_2(\psi(p))$ satisfies all the conditions we required.
\end{theorem}


\begin{definition}
    Suppose $M$ is a topological space, and let $\mathcal{A}=\left\{A_\alpha\right\}$ be an arbitrary open cover of $M$.
    A \textbf{partition of unity subordinate to $\mathcal{A}$} is an indexed family
    \begin{equation*}
        \left\{\rho_\alpha: \rho_\alpha: M \rightarrow \mathbb{R} \text{ is continue}\right\}
    \end{equation*}
    with the following properties:
    \begin{enumerate}[label=(\roman*)]
        \item $\rho_\alpha \prec A_\alpha$
              $(\supp \rho_\alpha \subset A_\alpha)$ for all $\alpha$
        \item  The family of supports $\left\{\operatorname{supp} \rho_\alpha\right\}$ is locally finite.
        \item  $\sum_{\alpha} \rho_\alpha(x) \equiv 1$ on $M$
    \end{enumerate}
\end{definition}

\begin{theorem}
    Suppose $M$ is a smooth manifold with or without boundary, and $\left\{U_\alpha\right\}_{\alpha \in A}$ is any indexed open cover of $M$. Then there exists a smooth partition of unity subordinate to $\left\{U_\alpha\right\}_{\alpha \in A}$.

    Proof.
    We can find nonnegative functions $\varphi_j \in C^{\infty}(M)$ so that $\overline{V}_j\prec\varphi_j \prec  W_j$ on $\overline{V}_j$ and $\operatorname{supp}\left(\varphi_j\right) \subset W_j$.
    Since $\mathcal{W}$ is a locally finite cover, $\varphi=\sum \varphi_j$ is a well-defined smooth function on $M$. Since each $\varphi_j$ is nonnegative, and $\mathcal{V}$ is a cover of $M, \varphi$ is strictly positive on $M$. It follows that the functions $\psi_j=\frac{\varphi_j}{\varphi}$ are smooth and satisfy $0 \leq \psi_j \leq 1$ and $\sum_j \psi_j=1$.

    Let
    \begin{equation*}
        \rho_\alpha=\sum_{W_j \subset U_\alpha} \psi_j
    \end{equation*}
    Note that the right hand side is a finite sum near each point, so it does define a smooth function. Clearly the family $\left\{\rho_\alpha\right\}$ is a partition of unity subordinate to $\left\{U_\alpha\right\}$.
\end{theorem}



\subsection{Application}

\begin{theorem}[Smooth Urysohn lemma]
    Let $M$ be a manifold, $F \subset M$ is a closed subset, and $U \subset M$ an open subset that contains $F$. Then there is a "bump" function $\varphi \in C^{\infty}(M)$ so that $0 \leq \varphi \leq 1, \varphi \equiv 1$ on $F$ and $\operatorname{supp}(\varphi) \subset U$.

    Proof. Let $\left\{\rho_1, \rho_2\right\}$ be a partition of unity subordinate to the open cover $\left\{U, M \backslash F\right\}$. Then $\varphi=\rho_1$ is what we need: $\rho_1=1$ on $F$ since $\rho_2=0$ on $F$.
\end{theorem}




\begin{theorem}[Whitney Approximation Theorem]
    Let $M$ be a smooth manifold, $F \subset M$ a closed subset and $k$ be a positive integer.
    Then for any continuous function $f: M \rightarrow \mathbb{R}^k$ which is smooth on $F$ and any positive continuous function $\delta: M \rightarrow \mathbb{R}_{>0}$, there exists $f \in C^{\infty}(M)$ so that
    \begin{equation*}
        f(p)=g(p), \quad \forall p \in F
    \end{equation*}
    and
    \begin{equation*}
        |f(p)-g(p)|<\delta(p), \quad \forall p \in M
    \end{equation*}

    Proof.
    By definition, there exists an open set $U \supset A$ and a smooth function $f_0$ defined on $U$ so that $f_0=f$ on $F$. Let
    \begin{equation*}
        U_0
        =
        \left\{p \in U:\left|f_0(p)-f(p)\right|<\delta(p)\right\}
    \end{equation*}
    Then $U_0$ is open in $M$ and $U_0 \supset F$.


    Next we construct a open cover of $M \backslash F$.
    For any $q \in M \backslash F$, we let
    \begin{equation*}
        U_q
        =
        \{p \in M \backslash A:|f(p)-f(q)|<\delta(p)\}
    \end{equation*}
    Then $\left\{U_q \mid q \in M \backslash F\right\}$ is an open covering of $M \backslash F$.
    Now let $\left\{\rho_0, \rho_q: q \in M\right\}$ be P.O.U. subordinate to the open cover $\left\{U_0, U_q: q \in M\right\}$ of $M$, and define a function on $M$ via
    \begin{equation*}
        g(p)
        =
        \rho_0(p) f_0(p)+\sum_{q \in M} \rho_q(p) f(q) .
    \end{equation*}
    Since the summation is locally finite, $g$ is smooth. Also by definition, $g=f_0=f$ on $F$. Moreover, for any $q \in M$ one has
    \begin{equation*}
        \begin{aligned}
            \left| g(p)-f(p)\right| & =\left|\rho_0(p) f_0(p)+\sum_q \rho_q(p) f(q)-\rho_0(p) f(p)-\sum_q \rho_q(p) f(p)\right| \\
                                    & \leq \rho_0(p)\left|f_0(p)-f(p)\right|+\sum_q \rho_q(p)|f(q)-f(p)|                        \\
                                    & <\rho_0(p) \delta(p)+\sum_q \rho_q(p) \delta(p)                                           \\
                                    & =\delta(p)
        \end{aligned}
    \end{equation*}
\end{theorem}
\begin{corollary}[Tietze]
    Let $M$ be a smooth manifold and closed set $F \subset$. If $f$ is smooth on $F$, then there exists $g \in C^\infty\left(M\right)$ that $f=g$ on $F$.
\end{corollary}

\begin{theorem}
    [Existence of Smooth Exhaustion Function]
    Every smooth manifold with or without boundary admits a smooth positive exhaustion function.
    \begin{proof}
        Let $M$ be a smooth manifold with or without boundary, let $\left\{V_j\right\}_{j=1}^{\infty}$ be any countable open cover of $M$ by precompact open subsets, and let $\left\{\psi_j\right\}$ be a smooth partition of unity subordinate to this cover. Define $f \in C^{\infty}(M)$ by
        \begin{equation*}
            f(p)=\sum_{j=1}^{\infty} j \psi_j(p)
        \end{equation*}
        Then $f$ is smooth because only finitely many terms are nonzero in a neighborhood of any point, and positive because $f(p) \geq \sum_j \psi_j(p)=1$.

        To see that $f$ is an exhaustion function, let $c \in \mathbb{R}$ be arbitrary, and choose a positive integer $N>c$. If $p \notin \bigcup_{j=1}^N \bar{V}_j$, then $\psi_j(p)=0$ for $1 \leq j \leq N$, so
        \begin{equation*}
            f(p)=\sum_{j=N+1}^{\infty} j \psi_j(p) \geq \sum_{j=N+1}^{\infty} N \psi_j(p)=N \sum_{j=1}^{\infty} \psi_j(p)=N>c .
        \end{equation*}
        Equivalently, if $f(p) \leq c$, then $p \in \bigcup_{j=1}^N \bar{V}_j$. Thus $f^{-1}((-\infty, c])$ is a closed subset of the compact set $\bigcup_{j=1}^N \bar{V}_j$ and is therefore compact.
    \end{proof}
\end{theorem}

\begin{theorem}[Level Sets of Smooth Functions]
    Let $M$ be a smooth manifold. If $K$ is any closed subset of $M$, there is a smooth nonnegative function $f: M \rightarrow \mathbb{R}$ such that $f^{-1}(0)=K$.
    \begin{proof}
        We begin with the special case in which $M=\mathbb{R}^n$ and $K \subseteq \mathbb{R}^n$ is a closed subset. For each $x \in M \backslash K$, there is a positive number $r \leq 1$ such that $B_r(x) \subseteq$ $M \backslash K$. By Proposition A.16, $M \backslash K$ is the union of countably many such balls $\left\{B_{r_i}\left(x_i\right)\right\}_{i=1}^{\infty}$.

        Let $h: \mathbb{R}^n \rightarrow \mathbb{R}$ be a smooth bump function that is equal to 1 on $\bar{B}_{1 / 2}(0)$ and supported in $B_1(0)$. For each positive integer $i$, let $C_i \geq 1$ be a constant that bounds the absolute values of $h$ and all of its partial derivatives up through order $i$. Define $f: \mathbb{R}^n \rightarrow \mathbb{R}$ by
        \begin{equation*}
            f(x)=\sum_{i=1}^{\infty} \frac{\left(r_i\right)^i}{2^i C_i} h\left(\frac{x-x_i}{r_i}\right)
        \end{equation*}

        The terms of the series are bounded in absolute value by those of the convergent series $\sum_i 1 / 2^i$, so the entire series converges uniformly to a continuous function by the Weierstrass $M$-test. Because the $i$ th term is positive exactly when $x \in B_{r_i}\left(x_i\right)$, it follows that $f$ is zero in $K$ and positive elsewhere.

        It remains only to show that $f$ is smooth. We have already shown that it is continuous, so suppose $k \geq 1$ and assume by induction that all partial derivatives of $f$ of order less than $k$ exist and are continuous. By the chain rule and induction, every $k$ th partial derivative of the $i$ th term in the series can be written in the form
        \begin{equation*}
            \frac{\left(r_i\right)^{i-k}}{2^i C_i} D_k h\left(\frac{x-x_i}{r_i}\right),
        \end{equation*}
        where $D_k h$ is some $k$ th partial derivative of $h$. By our choices of $r_i$ and $C_i$, as soon as $i \geq k$, each of these terms is bounded in absolute value by $1 / 2^i$, so the differentiated series also converges uniformly to a continuous function. It then follows from Theorem C. 31 that the $k$ th partial derivatives of $f$ exist and are continuous. This completes the induction, and shows that $f$ is smooth.

        Now let $M$ be an arbitrary smooth manifold, and $K \subseteq M$ be any closed subset. Let $\left\{B_\alpha\right\}$ be an open cover of $M$ by smooth coordinate balls, and let $\left\{\psi_\alpha\right\}$ be a subordinate partition of unity. Since each $B_\alpha$ is diffeomorphic to $\mathbb{R}^n$, the preceding argument shows that for each $\alpha$ there is a smooth nonnegative function $f_\alpha: B_\alpha \rightarrow \mathbb{R}$ such that $f_\alpha^{-1}(0)=B_\alpha \cap K$. The function $f=\sum_\alpha \psi_\alpha f_\alpha$ does the trick.
    \end{proof}
\end{theorem}


\subsection{In paracompact Hausdorff space}

\begin{lemma}[Shrinking lemma]
    Let $X$ be a paracompact Hausdorff space; let $\left\{U_\alpha\right\}_{\alpha \in J}$ be open cover of $X$.
    Then there exists a locally finite open cover $\left\{V_\alpha\right\}_{\alpha \in J}$ of  $X$ such that $\overline{V}_\alpha \subset U_\alpha$ for each $\alpha$.

    Proof. Let $\mathcal{A}$ be the collection of all open sets $A$ such that $\overline{A}$ is contained in some element of the collection $\left\{U_\alpha\right\}$.
    Regularity of $X$ implies that $\mathcal{A}$ covers $X$.
    Since $X$ is paracompact, we can find a locally finite collection $\mathcal{B}=\left\{B_\beta\right\}_{\beta \in K}$ of open sets covering $X$ that refines $\mathcal{A}$

    Let us index $\mathcal{B}$ bijectively with some index set $K$, then the general element of $\mathscr{B}$ can be denoted $B_\beta$, for $\beta \in K$, and $\left\{B_\beta\right\}_{\beta \in K}$ is a locally finite indexed famuly. Since $\mathscr{B}$ refines $\mathscr{A}$, we can define a function $f: K \rightarrow J$ by choosing, for each $\beta$ in $K$, an element $f(\beta) \in J$ such that
    \begin{equation*}
        \overline{B}_\beta \subset U_{f(\beta)}
    \end{equation*}
    Then for each $\alpha \in J$, we define $V_\alpha$ to be the union of many $B_\beta$
    \begin{equation*}
        V_\alpha
        =
        \bigcup_{f(\beta)=\alpha}B_\beta
    \end{equation*}
    Because the collection $\mathcal{B}_\alpha$ is locally finite, $\overline{V}_\alpha= \bigcup \overline{B}_\beta$, so that $\overline{V}_\alpha \subset U_\alpha$.

    Finally, we check local finiteness. Given $x \in X$, choose a neighborhood $W$ of $x$ such that $W$ intersects $B_\beta$ for only finitely many values of $\beta$, say $\beta=\beta_1, \ldots, \beta_K$. Then $W$ can intersect $V_\alpha$ only if $\alpha$ is one of the indices $f\left(\beta_1\right), \ldots, f\left(\beta_K\right)$.
\end{lemma}

\begin{theorem}
    Let $X$ be a paracompact Hausdorff space; let $\left\{U_\alpha\right\}_{\alpha \in J}$ be an indexed open covering of $X$. Then there exists a partition of unity on $X$ subordinate to $\left\{U_\alpha\right\}$.

    Proof.
    We begin by applying the shrinking lemma twice, to find locally finite indexed famles of open sets $\left\{W_\alpha\right\}$ and $\left\{V_\alpha\right\}$ covering $X$, such that
    \begin{equation*}
        W_\alpha \subset \overline{W}_\alpha \subset V_\alpha \subset \overline{V}_\alpha \subset U_\alpha
    \end{equation*}
    for each $\alpha$ Since $X$ is normal, we may choose, for each $\alpha$, a continuous function $\varphi_\alpha: X \rightarrow[0,1]$ such that $\varphi_\alpha\left(\overline{W}_\alpha\right)=\{1\}$ and $\varphi_\alpha\left(X-V_\alpha\right)=\{0\}$.
    Since $\varphi_\alpha$ is nonzero only at points of $V_\alpha$, we have
    \begin{equation*}
        \supp \varphi_\alpha\subset \overline{V}_\alpha \subset U_\alpha
    \end{equation*}
    Furthermore, the indexed family $\left\{\overline{V}_\alpha\right\}$ is locally finite (since an open set intersects $\overline{V}_\alpha$ only if it intersects $V_\alpha$ ); hence the indexed family  $\supp \varphi_\alpha$ is also locally finite. Note that because $\left\{W_\alpha\right\}$ covers $X$, for any given $x$ at least one of the functions $\psi_\alpha$ is positive at $x$.

    We can now make sense of the formally infinite sum
    \begin{equation*}
        \varphi(x)=\sum_\alpha \varphi_\alpha(x)
    \end{equation*}
    and define
    \begin{equation*}
        \rho_\alpha(x)
        =
        \frac{\varphi_\alpha(x)}{\varphi(x)}
    \end{equation*}
    to obtain our desired partition of unity.
\end{theorem}

\section{Sard's Theorem}
\begin{definition}
    We say that a subset $A \subseteq M$ \textbf{has measure zero in $M$} if for every smooth chart $(U, \varphi)$ for $M$, the subset $\varphi(A \cap U) \subseteq \mathbb{R}^n$ has $n$-dimensional measure zero.


\end{definition}

\begin{lemma}
    Let $M$ be a smooth $n$-manifold with or without boundary and $A \subseteq M$. Suppose that for some collection $\left\{\left(U_\alpha, \varphi_\alpha\right)\right\}$ of smooth charts whose domains cover $A, \varphi_\alpha\left(A \cap U_\alpha\right)$ has measure zero in $\mathbb{R}^n$ for each $\alpha$. Then $A$ has measure zero in $M$.
\end{lemma}

\begin{theorem}
    Suppose $M$ and $N$ are differential $m$-manifolds with or without boundary, $F: M \rightarrow N$ is a $C^1$ map, and $A \subseteq M$ is a subset of measure zero. Then $F(A)$ has measure zero in $N$.
\end{theorem}

\begin{corollary}
    Suppose $M$ and $N$ are differential manifolds with or without boundary, $\dim M \leq \dim N$, and $F \in C^1\left(M,N\right)$.
    If $A \subset M$ is a subset of measure zero, then $F(A)$ has measure zero in $N$.
\end{corollary}




\begin{definition}
    If $f: M \rightarrow N$ is a smooth map,

    (1) A point $p \in M$ is said to be a \textbf{regular point of $f$} if $d f_p: T_p M \rightarrow$ $T_{f(p)} N$ is surjective $\left(\rank_p f =\dim N\right)$.
    A point $c \in N$ is said to be a \textbf{regular value} of $f$ if every point of the level set $f^{-1}(c)$ is a regular point.

    (2) A point $q \in M$ is said to be a \textbf{critical point of $f$} if $d f_p: T_p M \rightarrow$ $T_{f(p)} N$ is not surjective $\left(\rank_p f <\dim N\right)$.
    A point $d \in N$ is said to be a \textbf{critical value} of $f$ if there exists a critical point in level set $f^{-1}(d)$.
\end{definition}

\begin{theorem}[Sard's Theorem]
    Suppose $M$ and $N$ are manifolds, and $f: M \rightarrow N$ is a smooth map.
    Then the set of critical values of $f$ has measure zero in $N$.
\end{theorem}


\begin{corollary}
    Suppose $M$ and $N$ are smooth manifolds with or without boundary, and $F: M \rightarrow N$ is a smooth map. If $\operatorname{dim} M<\operatorname{dim} N$, then $F(M)$ has measure zero in $N$.
\end{corollary}

\begin{corollary}
    Suppose $M$ is a smooth manifold with or without boundary, and $S \subseteq M$ is an immersed submanifold with or without boundary. If $\operatorname{dim} S<\operatorname{dim} M$, then $S$ has measure zero in $M$.
\end{corollary}







\chapter{Tangent and Cotangent Spaces}
\minitoc

\section{Basic Definitions}
\begin{definition}
    Let $M$ be a manifold, $p$ be a point in $M$.
    \begin{enumerate}
        \item
              The \textbf{germ} of a $C^\infty$ function at $p$ in $M$ is defined
              \begin{equation*}
                  C^\infty_p\left(M\right)
                  :=
                  \varinjlim_{U \ni p} C^{\infty}(U)
              \end{equation*}
        \item
              A \textbf{tangent vector} at $p$ is  a derivation of $\mathbb{R}$-algebra $ C^\infty_p\left(M\right)$
              ( a linear map $D: C_p^{\infty}(M) \rightarrow \mathbb{R}$ such that  $D(f g)=(D f) g(p)+f(p) D g$)
        \item
              The tangent vectors at $p$ form a vector space
              \begin{equation*}
                  T_pM
                  :=
                  \Der_{\mathbb{R}} \left(C_p^{\infty}(M), \mathbb{R}\right)
              \end{equation*}
              called the \textbf{tangent space} of $M$ at $p$.
        \item
              A \textbf{cotangent vector} at $p$ is a linear functional on the tangent space $T_p M$.
              And the \textbf{cotangent space} of $M$ at $p$ is the dual vector space
              \begin{equation*}
                  T_p^* M
                  :=
                  \Hom_{\mathbb{R}}\left(T_p M, \mathbb{R}\right)
              \end{equation*}
    \end{enumerate}
\end{definition}


\begin{definition}
    We define the \textbf{tangent bundle} of $M$:
    \begin{equation*}
        T M
        :=
        \coprod_{p \in M} T_p M
    \end{equation*}
    \begin{remark}
        We usually write an element of this disjoint union as an ordered pair $(p, X)$ of $X_p$, with $p \in M$ and $X \in T_p M$. The tangent bundle comes equipped with a \textbf{natural projection map}
        \begin{equation*}
            \pi: T M \rightarrow M
        \end{equation*}
        which sends each vector in $T_p M$ to the point $p$ at which it is tangent: $\pi(p, X)=p$.
    \end{remark}
\end{definition}

\begin{proposition}
    The tangent bundle $T M$ has a natural topology and smooth structure that make it into a
    $2 n$-dimensional smooth manifold.
    With respect to this structure, the projection $\pi$ is smooth.
\end{proposition}

\begin{definition}
    The cotan
\end{definition}







\section{The Differential of a Smooth Map}
% The Differential of a Map
\begin{definition}
    Let $f:M \rightarrow N$ be smooth map between manifolds and a point $p\in M$.
    We define a map induced by $f$ on the tangent spaces
    \begin{equation*}
        \mathrm{d} f_p: T_p M \rightarrow T_{f(p)} N
    \end{equation*}
    as follows.
    Given $v \in T_p M$, we let $\mathrm{d} f_p(v)$ be the derivation at $f(p)$ that acts on $C^{\infty}_{f(p)}(N)$ by the rule
    \begin{equation*}
        \left\langle \mathrm{d} f_p(v),  \varphi \right\rangle
        =
        \left\langle  v,  \varphi \circ f \right\rangle \quad \text{ for all } \varphi \in C_{f(p)}^\infty(N)
    \end{equation*}
\end{definition}

\begin{proposition}
    There is a smooth map called the \textbf{pushforward} of $f$
    \begin{equation*}
        f_*: T M \rightarrow T N ,\quad (p, v) \mapsto (f(p), \mathrm{d} f_p(v))
    \end{equation*}
\end{proposition}





\begin{proposition}
    Tangent functor $T: \mathbf{Man} \rightarrow \mathbf{VectBun}$ is covariant
\end{proposition}


\subsection{Computations in Coordinates} % Computations in Coordinates

Given a chart $\left(U, \varphi, x^i\right)$ covering a point $p \in M$, $\mathrm{d}\varphi_p$ is an isomorphism between $T_p M$ and $T_p\mathbb{R}^n$.
Thus we can choose a basis of $T_p M$ by
\begin{equation*}
    \left.\frac{\partial}{\partial x^i}\right|_p
    :=
    (\mathrm{d}\varphi_p)^{-1}\left(\left.\frac{\partial}{\partial x^i}\right|_{\hat{p}}\right)
\end{equation*}


\begin{proposition}
    Given a smooth map $f: M \rightarrow N$, and charts $\left(U, \varphi, x^i\right)$ and $\left(V, \psi, y^j\right)$ for $M$ and $N$ respectively, then
    \begin{equation*}
        df_p \left(\left.\frac{\partial}{\partial x^i}\right|_p\right)
        =
        \sum_j\frac{\partial y^j}{\partial x^i}\left(\hat{p}\right)\cdot\left.\frac{\partial}{\partial y^j}\right|_{f(p)}
    \end{equation*}
    where $(y^1, \ldots, y^n)=\psi \circ f\circ \varphi^{-1}(x^1, \ldots, x^n)$.
    Thus the representation matrix of $df_p$ with respect to the bases $\left\{\left.\frac{\partial}{\partial x^i}\right|_p\right\}$ and $\left\{\left.\frac{\partial}{\partial y^j}\right|_{f(p)}\right\}$ is the Jacobian matrix of $\tilde{f}$ at $\hat{p}$.
\end{proposition}

\begin{corollary}
    Suppose that $\left(U,x^i\right)$ and $(\tilde{U},\tilde{x}^i)$ are smooth charts for $M$ centered at $p$, then
    \begin{equation*}
        df_p \left(\left.\frac{\partial}{\partial x^i}\right|_p\right)
        =
        \frac{\partial \tilde{x}^j}{\partial x^i}(\hat{p})\cdot \left.\frac{\partial}{\partial \tilde{x}^j}\right|_p
    \end{equation*}
\end{corollary}


\section{Vector fields}
\begin{definition}
    A \textbf{vector field} on $M$ is a section of $\pi: T M \rightarrow M$, i.e. a map $X: M \rightarrow T M$ such that $\pi \circ X=\operatorname{id}_M$.
    The \textbf{support} of a vector field $X$ is the closure of the set $\left\{p \in M: X(p) \neq 0\right\}$.
    We say that $X$ is a \textbf{smooth vector field} if it is smooth as a map from $M$ to $T M$.
    \begin{remark}
        If $\left(U,x^i\right)$ is a smooth chart, we can write $X$ (omit the base point) in coordinates locally as
        \begin{equation*}
            X(p)
            =
            \sum_i X^i(p) \cdot \left.\frac{\partial}{\partial x^i}\right|_p
        \end{equation*}
        This $n$ functions $X^i: U \rightarrow \mathbb{R}$ are called the \textbf{components funtions} of $X$ in the given chart.
    \end{remark}
\end{definition}

\begin{proposition}[Smoothness Criterion for Vector Fields]
    A vector field $X$ is smooth if and only if its component functions are smooth.
\end{proposition}

\begin{proposition}
    The set of all smooth vector fields on $M$ is a $C^\infty(M)$-module:
    \begin{equation*}
        \mathfrak{X}(M) =\Der_{\mathbb{R}}\left(C^\infty(M),C^\infty(M)\right)
    \end{equation*}
\end{proposition}



\begin{definition}
    Let $f: M \rightarrow N$ be a smooth map, and let $X\in \mathfrak{X}(M)$.
    The vector field $X$ and $Y$ are said to be \textbf{$f$-related} if for every $p \in M$,
    \begin{equation*}
        f_{*,p}(X_p)=Y_{f(p)}
    \end{equation*}
\end{definition}

\begin{proposition}
    If $f$ is a diffeomorphism, then

    \begin{enumerate}
        \item
              for every $X \in \mathfrak{X}(M)$, there is a unique vector field $Y$ on $N$ that is $f$-related to $X$ called the \textbf{pushforward} of $X$ under $f$.
              Explicitly,
              \begin{equation*}
                  (f_{*}X)_q=f_{*,f^{-1}(q)}(X_{f^{-1}(q)})
              \end{equation*}
    \end{enumerate}
\end{proposition}


\subsection{Integral Curves}
Given a smooth curve $\gamma: J \rightarrow M$ and $t_0 \in J$, we define the \textbf{velocity} of $\gamma$ at $t_{0}$, denoted by $\gamma^{\prime}\left(t_0\right)$, to be the vector
\begin{equation*}
    \gamma^{\prime}\left(t_0\right)
    :=
    \gamma_{*,t_0}\left(\left.\frac{d}{d t}\right|_{t_0}\right) \in T_{\gamma\left(t_0\right)} M
\end{equation*}
where $d /\left.d t\right|_{t_0}$ is the standard coordinate basis vector in $T_{t_0} \mathbb{R}$.

\begin{definition}
    Let $X\in \mathfrak{X}(M)$ be a smooth vector field.
    The curve $\gamma: J \rightarrow M$ is called an \textbf{integral curve} of $X$ if $\gamma^{\prime}(t)=X_{\gamma(t)}$ for every $t \in J$
\end{definition}

\subsection{Flow}

\begin{definition}
    A \textbf{global flow} on $M$ is a left $\mathbb{R}$-action $\theta: \mathbb{R} \times M \rightarrow M$.
\end{definition}








\begin{definition}
    A \textbf{flow} on $M$ is a continuous map $\theta:\mathcal{D} \rightarrow M$
    \begin{enumerate}[label=(\roman*)]
        \item
              $\mathcal{D}\subseteq \mathbb{R} \times M$ is the \textbf{flow domain} for $M$, that is, $\mathcal{D}^{(p)}:=\left\{t\in \mathbb{R}|(p,t)\in \mathcal{D}\right\}$ is an open interval containing $0$ for each $p \in M$
        \item
              $\theta(0, p)=p$ for all $p \in M$, and
        \item
              for all $s\in \mathcal{D}^{(p)}$ and $t\in \mathcal{D}^{\theta(s,p)}$ such that $s+t\in \mathcal{D}^{(p)}$, we have $\theta\left(s, \theta(t, p)\right)=\theta(s+t, p)$

    \end{enumerate}
    \begin{remark}
        For each $p\in M$, define a curve $\theta^{(p)}: \mathcal{D}^{(p)} \rightarrow M$ by $\theta^{(p)}(t)=\theta(t, p)$
    \end{remark}
\end{definition}

\begin{definition}
    Given a smooth flow $\theta$ on $M$, the \textbf{infinitesimal generator} of $\theta$ is the vector field $X$ defined by
    \begin{equation*}
        p\mapsto \theta^{(p)\prime}(0)
    \end{equation*}
\end{definition}

\begin{proposition}
    The infinitesimal generator $X$ of $\theta$ is smooth, and
    each curve $\theta^{(p)}$ is an integral curve of $X$.
\end{proposition}


\begin{theorem}[Fundamental Theorem on Flows]

\end{theorem}








\chapter{Vector Bundles}

\section{Vector Bundles}

\begin{definition}
    Let $M$ be a topological space. A \textbf{vector bundle of rank $k$ over $M$} is a topological space $E$ together with a continuous surjective map
    \[
        \pi:E\rightarrow M
    \]
    such that
    \begin{enumerate}[label=(\roman*)]
        \item for each $p\in M$, the fiber $E_p=\pi^{-1}(p)$ is a $k$-dimensional real vector space;
        \item every $p\in M$ has an open neighborhood $U$ and a diffeomorphism
              \[
                  \Phi:\pi^{-1}(U)\rightarrow U\times \mathbb{R}^k
              \]
              satisfying $\pi_1\circ\Phi=\pi$, where $\pi_1:U\times\mathbb{R}^k\rightarrow U$ is projection;
        \item for each $q\in U$, the restricted map
              \[
                  \Phi|_{E_q}:E_q\rightarrow \{q\}\times \mathbb{R}^k\cong \mathbb{R}^k
              \]
              is a linear isomorphism.
    \end{enumerate}
    The manifold $E$ is the \textbf{total space}, $M$ is the \textbf{base}, $\pi$ is the \textbf{projection}, and $\Phi$ is a \textbf{smooth local trivialization}.
\end{definition}



\begin{example}
    The tangent bundle $TM\rightarrow M$ is a rank-$n$ vector bundle when $\dim M=n$. If $(U,x^1,\ldots,x^n)$ is a smooth chart, then
    \[
        \Phi\left(v^i\frac{\partial}{\partial x^i}\bigg|_p\right)
        =
        (p,(v^1,\ldots,v^n))
    \]
    is a local trivialization of $TM$ over $U$.
\end{example}

\begin{definition}
    If $\Phi:\pi^{-1}(U)\rightarrow U\times\mathbb{R}^k$ and $\Psi:\pi^{-1}(V)\rightarrow V\times\mathbb{R}^k$ are local trivializations of a rank-$k$ vector bundle, then on $U\cap V$ their \textbf{transition function} is the smooth map
    \[
        \tau:U\cap V\rightarrow \GL(k,\mathbb{R})
    \]
    defined by
    \[
        \Psi\circ\Phi^{-1}(p,v)=(p,\tau(p)v).
    \]
    Thus changes of vector bundle coordinates are linear on each fiber, and the corresponding linear maps vary smoothly with the base point.
\end{definition}

\begin{lemma}[Vector Bundle Chart Lemma]
    Let $M$ be a smooth manifold and let $\pi:E\rightarrow M$ be a surjective map whose fibers are $k$-dimensional vector spaces. Suppose there is an open cover $\{U_\alpha\}$ of $M$ and bijections
    \[
        \Phi_\alpha:\pi^{-1}(U_\alpha)\rightarrow U_\alpha\times\mathbb{R}^k
    \]
    such that $\pi_1\circ\Phi_\alpha=\pi$, each $\Phi_\alpha|_{E_p}$ is a linear isomorphism, and the transition maps have the form
    \[
        \Phi_\beta\circ\Phi_\alpha^{-1}(p,v)=(p,\tau_{\beta\alpha}(p)v)
    \]
    with $\tau_{\beta\alpha}:U_\alpha\cap U_\beta\rightarrow \GL(k,\mathbb{R})$ smooth. Then $E$ has a unique smooth manifold structure for which $\pi:E\rightarrow M$ is a smooth vector bundle and the $\Phi_\alpha$ are smooth local trivializations.
\end{lemma}

\begin{example}[Whitney Sum]
    If $\pi':E'\rightarrow M$ and $\pi'':E''\rightarrow M$ are smooth vector bundles of ranks $k'$ and $k''$, their \textbf{Whitney sum} is
    \[
        E'\oplus E''
        =
        \coprod_{p\in M}\left(E'_p\oplus E''_p\right).
    \]
    Using local trivializations of $E'$ and $E''$ over the same open set gives local trivializations
    \[
        (E'\oplus E'')|_U\cong U\times\mathbb{R}^{k'+k''}.
    \]
    Hence $E'\oplus E''$ is a smooth vector bundle of rank $k'+k''$.
\end{example}

\begin{example}[Restriction]
    If $\pi:E\rightarrow M$ is a smooth vector bundle and $S\subseteq M$ is an embedded submanifold, then
    \[
        E|_S=\pi^{-1}(S)
    \]
    is a smooth vector bundle over $S$ with projection $\pi|_{E|_S}:E|_S\rightarrow S$. Its fiber over $p\in S$ is still $E_p$.
\end{example}

\section{Local and Global Sections}

\begin{definition}
    Let $\pi:E\rightarrow M$ be a vector bundle.
    \begin{enumerate}
        \item
              A \textbf{global section} of $E$ is a continuous map
              $\sigma:M\rightarrow E$
              such that $\pi\circ\sigma=\id_M$.

              A \textbf{local section} of $E$ is a continuous map
              $\sigma:U\rightarrow E$
              such that $\pi\circ\sigma=\id_U$.

        \item

    \end{enumerate}
\end{definition}



\begin{definition}
    Let $E\rightarrow M$ be a vector bundle and a open subset $U\subseteq M$. A \textbf{local frame} for $E$ over $U$ is a $k$-tuple of local sections
    \[
        (\sigma_1,\ldots,\sigma_k)
    \]
    such that $\left(\sigma_1(p),\ldots,\sigma_k(p)\right)$ is a basis for $E_p$ for every $p\in U$. A frame over all of $M$ is a \textbf{global frame}.
\end{definition}

\begin{proposition}[Frames and Trivializations]
    Smooth local trivializations and smooth local frames determine each other.
    If $\Phi:\pi^{-1}(U)\rightarrow U\times\mathbb{R}^k$ is a local trivialization and $(e_1,\ldots,e_k)$ is the standard basis of $\mathbb{R}^k$, then
    \[
        \sigma_i(p)=\Phi^{-1}(p,e_i)
    \]
    is a smooth local frame over $U$. Conversely, if $(\sigma_1,\ldots,\sigma_k)$ is a smooth local frame, then
    \[
        \Phi\left(a^i\sigma_i(p)\right)=(p,(a^1,\ldots,a^k))
    \]
    defines a smooth local trivialization.
\end{proposition}

\begin{corollary}
    A smooth rank-$k$ vector bundle is smoothly trivial if and only if it admits a smooth global frame.
\end{corollary}

\begin{proposition}[Completion of Local Frames]
    Let $E\rightarrow M$ be a smooth vector bundle of rank $k$. If smooth local sections $\sigma_1,\ldots,\sigma_j$ are linearly independent at a point $p$, then after shrinking to a neighborhood $U$ of $p$ they can be completed to a smooth local frame
    \[
        (\sigma_1,\ldots,\sigma_j,\sigma_{j+1},\ldots,\sigma_k)
    \]
    for $E$ over $U$.
\end{proposition}

\begin{corollary}
    Given $p\in M$ and $v\in E_p$, there is a smooth local section $\sigma$ of $E$ defined near $p$ such that $\sigma(p)=v$. Using a bump function, one may choose $\sigma$ to extend to a smooth global section whose support lies in a prescribed coordinate neighborhood.
\end{corollary}

\begin{proposition}[Local Frame Criterion for Smoothness]
    Let $(\sigma_1,\ldots,\sigma_k)$ be a smooth local frame for $E$ over $U$. A local section $\sigma:U\rightarrow E$ is smooth if and only if the component functions $a^1,\ldots,a^k$ in
    \[
        \sigma=a^i\sigma_i
    \]
    are smooth functions on $U$.
\end{proposition}

\section{Bundle Homomorphisms}

\begin{definition}
    Let $\pi:E\rightarrow M$ and $\pi':E'\rightarrow M'$ be smooth vector bundles. A \textbf{bundle homomorphism} from $E$ to $E'$ is a pair of smooth maps
    \[
        F:E\rightarrow E',
        \qquad
        f:M\rightarrow M'
    \]
    such that
    \[
        \pi'\circ F=f\circ\pi
    \]
    and each restricted map
    \[
        F|_{E_p}:E_p\rightarrow E'_{f(p)}
    \]
    is linear. If $M=M'$ and $f=\id_M$, then $F$ is said to be a bundle homomorphism \textbf{over $M$}.
\end{definition}

\begin{definition}
    A bundle homomorphism $F:E\rightarrow E'$ over $M$ is a \textbf{bundle isomorphism} if it is bijective and its inverse is also a smooth bundle homomorphism over $M$. Two vector bundles over $M$ are \textbf{smoothly isomorphic} if there exists such an isomorphism between them.
\end{definition}

\begin{proposition}
    A bijective smooth bundle homomorphism over $M$ is a smooth bundle isomorphism.
\end{proposition}

\begin{proposition}[Matrix Criterion]
    Let $E,E'\rightarrow M$ be smooth vector bundles and let $F:E\rightarrow E'$ be a bundle homomorphism over $M$. If $(\sigma_1,\ldots,\sigma_k)$ is a smooth local frame for $E$ and $(\tau_1,\ldots,\tau_\ell)$ is a smooth local frame for $E'$ over $U$, then
    \[
        F(\sigma_j)=F^i{}_j\,\tau_i.
    \]
    The map $F$ is smooth if and only if all coefficient functions $F^i{}_j$ are smooth on $U$ for every such choice of local frames.
\end{proposition}

\begin{lemma}[Characterization by Sections]
    Let $E,E'\rightarrow M$ be smooth vector bundles. A smooth bundle homomorphism $F:E\rightarrow E'$ over $M$ induces a $C^\infty(M)$-linear map
    \[
        \Gamma(E)\rightarrow \Gamma(E'),
        \qquad
        \sigma\mapsto F\circ\sigma.
    \]
    Conversely, every $C^\infty(M)$-linear map $\Gamma(E)\rightarrow \Gamma(E')$ is induced by a unique smooth bundle homomorphism over $M$.
\end{lemma}

\begin{example}
    If $f\in C^\infty(M)$, scalar multiplication
    \[
        v\mapsto f(\pi(v))v
    \]
    is a smooth bundle homomorphism $E\rightarrow E$ over $M$. On sections it is the map $\sigma\mapsto f\sigma$.
\end{example}

\begin{example}
    If $F:M\rightarrow N$ is smooth, then the differential
    \[
        dF:TM\rightarrow TN
    \]
    is a smooth bundle homomorphism covering $F$. On each fiber it is the linear map
    \[
        dF_p:T_pM\rightarrow T_{F(p)}N.
    \]
\end{example}

\section{Subbundles}

\begin{definition}
    Let $\pi:E\rightarrow M$ be a smooth vector bundle. A subset $D\subseteq E$ is a \textbf{smooth subbundle of rank $r$} if $D$ is an embedded submanifold of $E$, the restriction $\pi|_D:D\rightarrow M$ is a smooth vector bundle of rank $r$, and each fiber
    \[
        D_p=D\cap E_p
    \]
    is a vector subspace of $E_p$.
\end{definition}

\begin{lemma}[Local Frame Criterion for Subbundles]
    Let $D\subseteq E$ be a subset such that each $D_p=D\cap E_p$ is an $r$-dimensional vector subspace of $E_p$. Then $D$ is a smooth subbundle of $E$ if and only if every $p\in M$ has a neighborhood $U$ on which there exist smooth local sections $\sigma_1,\ldots,\sigma_r$ of $E$ satisfying
    \[
        D_q=\operatorname{span}\left(\sigma_1(q),\ldots,\sigma_r(q)\right)
    \]
    for all $q\in U$.
\end{lemma}

\begin{example}
    If $S\subseteq M$ is an embedded submanifold, then $TS$ is naturally a smooth subbundle of $TM|_S$. At each point $p\in S$,
    \[
        T_pS\subseteq T_pM.
    \]
\end{example}

\begin{theorem}[Kernel and Image Subbundles]
    Let $F:E\rightarrow E'$ be a smooth bundle homomorphism over $M$. If the rank of the linear map
    \[
        F_p:E_p\rightarrow E'_p
    \]
    is constant on $M$, then
    \[
        \Ker F=\coprod_{p\in M}\Ker F_p
        \qquad\text{and}\qquad
        \operatorname{Im}F=\coprod_{p\in M}\operatorname{Im}F_p
    \]
    are smooth subbundles of $E$ and $E'$ respectively.
\end{theorem}

\begin{definition}
    If $M\subseteq \mathbb{R}^n$ is an embedded submanifold, its \textbf{normal bundle} is
    \[
        NM
        =
        \coprod_{p\in M}N_pM,
        \qquad
        N_pM=\left(T_pM\right)^\perp\subseteq T_p\mathbb{R}^n\cong \mathbb{R}^n.
    \]
\end{definition}

\begin{corollary}
    If $M\subseteq \mathbb{R}^n$ is an embedded submanifold, then $NM$ is a smooth subbundle of $T\mathbb{R}^n|_M$. Moreover
    \[
        T\mathbb{R}^n|_M
        =
        TM\oplus NM.
    \]
\end{corollary}

\begin{remark}
    The central principle of vector bundles is local triviality: near each point, a vector bundle looks like $U\times\mathbb{R}^k$, but the transition functions can glue these local products nontrivially. Sections generalize vector fields, frames are the bundle analogue of bases, and bundle homomorphisms are fiberwise linear maps varying smoothly with the base point.
\end{remark}












\chapter{Tensors}
\section{}
Given a vector space $V$, denote by $T^k(V)$ the $k$-th tensor power $V^{\otimes k}$.






\begin{definition}
    Given a $k$-covector $\omega \in T^k(V^*)$ , we define
    \begin{enumerate}
        \item the \textbf{symmetrization} of $\omega$ to be the tensor
              \begin{equation*}
                  \operatorname{Sym}(\omega)(v_1, \ldots, v_k)
                  :=
                  \frac{1}{k!} \sum_{\sigma \in S_k} \omega\left(v_{\sigma(1)}, \ldots, v_{\sigma(k)}\right)
              \end{equation*}
        \item the \textbf{antisymmetrization} of $\omega$ to be the tensor
              \begin{equation*}
                  \operatorname{Alt}(\omega)(v_1, \ldots, v_k)
                  :=
                  \frac{1}{k!} \sum_{\sigma \in S_k} \operatorname{sgn}(\sigma) \omega\left(v_{\sigma(1)}, \ldots, v_{\sigma(k)}\right)
              \end{equation*}
    \end{enumerate}
\end{definition}

\begin{definition}
    Given $\omega \in T^k(V^*)$ and $\eta  \in T^l(V^*)$.
    The \textbf{wedge product} of $\omega $ and $\eta  $ is defined to be
    \begin{equation*}
        \omega \wedge \eta
        :=
        \frac{(k+l)!}{k!l!} \operatorname{Alt}(\omega \otimes \eta)
    \end{equation*}
    The \textbf{symmetric product} of $\omega $ and $\eta $ is defined to be
    \begin{equation*}
        \omega \odot \eta
        :=
        \frac{(k+l)!}{k!l!} \operatorname{Sym}(\omega \otimes \eta)
    \end{equation*}
\end{definition}




\begin{proposition}
    The wedge product has the following properties:
    \begin{enumerate}
        \item associative
        \item bilinear
        \item  Anti-commutative relation: $\omega \wedge \eta=(-1)^{k l} \eta \wedge \omega$
        \item For any $f^i \in V^*$ and $v^i \in V$, we have
              \begin{equation*}
                  f^1 \wedge \cdots \wedge f^k\left(v_1, \ldots, v_k\right)
                  =
                  \operatorname{det}
                  \left(f^i(v_i)\right)
              \end{equation*}
    \end{enumerate}
\end{proposition}

\begin{definition}
    Let $V$ be a vector space. A $(k,l)$-tensor $T$ on $V$ is a element of the vector space
    \begin{equation*}
        T^{k,l}(V)
        :=
        T^k(V) \otimes T^l\left(V^*\right)
    \end{equation*}

    \begin{remark}
        A multilinear map $T: \underbrace{V^* \times \cdots \times V^*}_{k} \times \underbrace{V \times \cdots \times V}_{l} \rightarrow \mathbb{R}$.
    \end{remark}
\end{definition}






\section{Pullbacks of Tensor Fields}

\begin{definition}
    Suppose $f: M \rightarrow N$ is a smooth map and $B$ is $k$-covariant tensor field on $N$.
    The \textbf{pullback} of $B$ by $f$ is the $k$-covariant tensor field $f^* B$ on $M$ defined pointwise by
    \begin{equation*}
        (f^* B)_p\left(X_1, \ldots, X_k\right)
        :=
        B_{f(p)}\left(f_{*,p}X_1, \ldots, f_{*,p}X_k\right)
    \end{equation*}
    for all $p \in M$ and $X_1, \ldots, X_k \in T_p M$.
\end{definition}






\begin{proposition}
    \begin{equation*}
        f^*(B_{i_1\ldots i_k} dy^{i_1} \otimes \cdots \otimes dy^{i_k})
        =
        B_{i_1\ldots i_k}\circ f \cdot d(y^{i_1} \circ f) \otimes \cdots \otimes d(y^{i_k} \circ f)
    \end{equation*}
\end{proposition}

\begin{proposition}
    The pullback is a contravariant functor from the category of smooth manifolds to the category of graded algebras:
    \begin{equation*}
        \begin{tikzcd}
            M \arrow[dd,"f"']& \mapsto &  T^{(*,0)}M\\
            &\mapsto& \\
            N & \mapsto & T^{(*,0)}N \arrow[uu,"f^*"']
        \end{tikzcd}
    \end{equation*}
    where $f^*$ is a graded algebra homomorphism, i.e.
    \begin{enumerate}
        \item
              $f^*(A\otimes B)=f^*A\otimes f^*B$
        \item $f^*(A+B )=f^*A + f^*B$ and $f^*(\lambda A)=\lambda f^*A$ for all $\lambda \in \mathbb{R}$.
    \end{enumerate}
\end{proposition}

\begin{corollary}
    \begin{enumerate}
        \item $F^*(fB)=f\circ F \cdot F^*B$ 
        \item The pullbak also commutes with the symmetrization, antisymmetrization, wedge product and symmetric product operations.
    \end{enumerate}
\end{corollary}


\begin{definition}
    The Lie derivative of a tensor field $A$ with respect to a vector field $X$ is defined by
    \begin{equation*}
        \mathcal{L}_X A
        :=
        \lim_{t \rightarrow 0} \frac{\left(\theta_{t}\right)^*_p A_{\theta_t(p)}-A_p}{t}
    \end{equation*}
\end{definition}








\chapter{Differential Forms}
\section{The Algebra of Alternating Tensors}
\begin{proposition}
    If $\omega\in \Lambda^n(V^*)$ and $T:V\rightarrow V$ is linear, then
    \[
        \omega(Tv_1,\ldots,Tv_n)
        =
        (\det T)\omega(v_1,\ldots,v_n).
    \]
    Thus top-degree alternating tensors transform by determinants.
\end{proposition}

\begin{definition}
    For $X\in V$, the \textbf{interior multiplication} or \textbf{contraction} by $X$ is the map
    \[
        \iota_X:\Lambda^k(V^*)\rightarrow \Lambda^{k-1}(V^*)
    \]
    defined by
    \[
        (\iota_X\omega)(X_1,\ldots,X_{k-1})
        =
        \omega(X,X_1,\ldots,X_{k-1}).
    \]
\end{definition}

\begin{proposition}
    Interior multiplication is an antiderivation of degree $-1$:
    \[
        \iota_X(\omega\wedge\eta)
        =
        (\iota_X\omega)\wedge\eta
        +
        (-1)^k\omega\wedge(\iota_X\eta),
        \qquad \omega\in \Lambda^k(V^*).
    \]
\end{proposition}

\section{Differential Forms on Manifolds}

\begin{definition}
    Let $M$ be a smooth manifold, possibly with boundary.
    \begin{enumerate}
        \item
              The bundle of alternating $k$-covectors is
              \[
                  \Lambda^kT^*M
                  =
                  \coprod_{p\in M}\Lambda^k(T_p^*M).
              \]
        \item
              A \textbf{differential $k$-form} on $M$ is a smooth section of $\Lambda^kT^*M$. The space of all smooth $k$-forms is denoted by
              \[
                  \Omega^k(M)=\Gamma(\Lambda^kT^*M).
              \]
              By convention, $\Omega^0(M)=C^\infty(M)$.
    \end{enumerate}

\end{definition}

\begin{definition}
    If $\omega\in\Omega^k(M)$ and $\eta\in\Omega^l(M)$, their wedge product is defined pointwise:
    \[
        (\omega\wedge\eta)_p
        =
        \omega_p\wedge\eta_p.
    \]
    Thus $\omega\wedge\eta\in\Omega^{k+l}(M)$.
\end{definition}

\begin{proposition}
    In a smooth chart $(U,x^1,\ldots,x^n)$, every $k$-form $\omega$ can be written uniquely as
    \[
        \omega
        =
        \sum_{i_1<\cdots<i_k}
        \omega_{i_1\cdots i_k}
        dx^{i_1}\wedge\cdots\wedge dx^{i_k},
    \]
    where the coefficient functions $\omega_{i_1\cdots i_k}$ are smooth on $U$.
\end{proposition}

\begin{example}
    On $\mathbb{R}^3$ with coordinates $(x,y,z)$,
    \[
        \Omega^1(\mathbb{R}^3)
        =
        \left\{
        P\,dx+Q\,dy+R\,dz
        \right\},
    \]
    and
    \[
        \Omega^2(\mathbb{R}^3)
        =
        \left\{
        A\,dy\wedge dz+B\,dz\wedge dx+C\,dx\wedge dy
        \right\}.
    \]
    The signs in the second expression are chosen to match the usual orientation conventions from vector calculus.
\end{example}

\begin{definition}
    Let $F:M\rightarrow N$ be a smooth map. The \textbf{pullback} of a $k$-form $\omega\in\Omega^k(N)$ is the $k$-form $F^*\omega\in\Omega^k(M)$ given by
    \[
        (F^*\omega)_p(v_1,\ldots,v_k)
        =
        \omega_{F(p)}(dF_p(v_1),\ldots,dF_p(v_k)).
    \]
\end{definition}

\begin{proposition}
    Pullbacks of forms satisfy
    \[
        F^*(\omega\wedge\eta)
        =
        F^*\omega\wedge F^*\eta,
        \qquad
        (G\circ F)^*
        =
        F^*\circ G^*,
        \qquad
        \id_M^*
        =
        \id.
    \]
    For a smooth function $f$ on $N$, this agrees with ordinary composition:
    \[
        F^*f=f\circ F.
    \]
\end{proposition}

\begin{proposition}
    Let $F:M\rightarrow N$ be smooth, and let $(y^1,\ldots,y^n)$ be local coordinates on $N$. If
    \[
        \omega
        =
        \sum_I \omega_I\,dy^{i_1}\wedge\cdots\wedge dy^{i_k},
    \]
    then
    \[
        F^*\omega
        =
        \sum_I
        (\omega_I\circ F)\,
        d(y^{i_1}\circ F)\wedge\cdots\wedge d(y^{i_k}\circ F).
    \]
\end{proposition}

\begin{definition}
    If $X$ is a smooth vector field on $M$ and $\omega\in\Omega^k(M)$, the interior product
    \[
        \iota_X\omega\in\Omega^{k-1}(M)
    \]
    is defined by
    \[
        (\iota_X\omega)_p=\iota_{X_p}(\omega_p).
    \]
\end{definition}

\section{Exterior Derivatives}

\begin{definition}
    The \textbf{exterior derivative} is the unique family of maps
    \[
        d:\Omega^k(M)\rightarrow \Omega^{k+1}(M)
    \]
    characterized by the following properties:
    \begin{enumerate}[label=(\arabic*)]
        \item if $f\in C^\infty(M)$, then $df$ is the ordinary differential of $f$;
        \item $d$ is $\mathbb{R}$-linear;
        \item $d(\omega\wedge\eta)=d\omega\wedge\eta+(-1)^k\omega\wedge d\eta$ for $\omega\in\Omega^k(M)$;
        \item $d(d\omega)=0$ for every differential form $\omega$.
    \end{enumerate}
\end{definition}

\begin{proposition}
    In local coordinates $(x^1,\ldots,x^n)$, if
    \[
        \omega
        =
        \sum_I \omega_I\,dx^{i_1}\wedge\cdots\wedge dx^{i_k},
    \]
    then
    \[
        d\omega
        =
        \sum_I d\omega_I\wedge dx^{i_1}\wedge\cdots\wedge dx^{i_k}
        =
        \sum_{j,I}
        \frac{\partial \omega_I}{\partial x^j}
        dx^j\wedge dx^{i_1}\wedge\cdots\wedge dx^{i_k}.
    \]
\end{proposition}




\begin{definition}
    A form $\omega$ is called \textbf{closed} if $d\omega=0$. It is called \textbf{exact} if there exists a form $\eta$ such that
    \[
        \omega=d\eta.
    \]
    Since $d^2=0$, every exact form is closed.
\end{definition}

\begin{proposition}
    Exterior differentiation commutes with pullback:
    \[
        F^*(d\omega)=d(F^*\omega)
    \]
    for every smooth map $F:M\rightarrow N$ and every $\omega\in\Omega^k(N)$.
\end{proposition}

\begin{proposition}[Invariant Formula for $d$]
    Let $\omega\in\Omega^k(M)$ and let $X_0,\ldots,X_k$ be smooth vector fields. Then
    \begin{align*}
        d\omega(X_0,\ldots,X_k)
         & =
        \sum_{i=0}^k
        (-1)^i
        X_i\bigl(\omega(X_0,\ldots,\widehat{X_i},\ldots,X_k)\bigr) \\
         & \quad+
        \sum_{0\leq i<j\leq k}
        (-1)^{i+j}
        \omega([X_i,X_j],X_0,\ldots,\widehat{X_i},\ldots,\widehat{X_j},\ldots,X_k).
    \end{align*}
    Here a hat means that the indicated argument is omitted.
\end{proposition}

\begin{corollary}
    For a $1$-form $\omega$ and vector fields $X,Y$,
    \[
        d\omega(X,Y)
        =
        X(\omega(Y))-Y(\omega(X))-\omega([X,Y]).
    \]
\end{corollary}



\subsubsection{Lie derivation of forms}




\begin{theorem}[Cartan's Magic Formula]
    On a smooth manifold, for every smooth vector field $V$ and every differential form $\omega$,
    \[
        \mathcal{L}_V\omega
        =
        \iota_V(d\omega)+d(\iota_V\omega).
    \]
\end{theorem}

\begin{corollary}
    The Lie derivative commutes with the exterior derivative:
    \[
        \mathcal{L}_V(d\omega)=d(\mathcal{L}_V\omega).
    \]
\end{corollary}



















\chapter{Submersions, Immersions, and Embeddings}
In this chapter, manifolds are assumed to be smooth manifolds with or without boundary unless otherwise stated.


\section{}
\begin{theorem}[Rank Theorem]
    Suppose $M$ and $N$ are smooth manifolds of dimensions $m$ and $n$, respectively, and
    $f : M \to N$
    is a smooth map with constant rank $r$. For each $p \in M$ there exist smooth charts
    $(U,\varphi)$
    for $M$ centered at $p$ and
    $(V,\psi)$
    for $N$ centered at $f(p)$ such that $f(U)\subseteq V$, in which $f$ has a coordinate representation of the form
    \[
        \hat{f}(x^1,\ldots,x^r,x^{r+1},\ldots,x^m)
        =
        (x^1,\ldots,x^r,0,\ldots,0).
    \]
\end{theorem}




\section{Submersions}


\begin{theorem}[Local Section Theorem]
    The smooth map $\pi : M \to N$ is a submersion if and only if every point $p\in M$ is in the image of a smooth local section of $\pi$.
\end{theorem}





\section{Embeddings}

\begin{lemma}[Transport of structure]
    Let $k\in \mathbb{N}$ and $f:M\to N$ be a homeomorphism between $C^k$-manifold $M$ and a topological space $N$.
    Then there is a unique $C^k$-structure on $N$ such that $f$ is a $C^k$-diffeomorphism.
    \begin{remark}
        If $f:M\to N$ is a homeomorphism between manifolds, $f$ is not necessarily a diffeomorphism. For example, the map $f:\mathbb{R}\to\mathbb{R}$ defined by $f(x)=x^3$ is a homeomorphism but not a diffeomorphism.
    \end{remark}
\end{lemma}

\begin{definition}
    A smooth map $f:M\rightarrow N$ is a \textbf{smooth embedding} if it is an immersion and a topological embedding. In this case, we say that $M$ is \textbf{smoothly embedded} in $N$, and $S=f(M) \cong M $ is an \textbf{embedded submanifold} of $N$.
\end{definition}
\begin{remark}
    By the transport of structure lemma, there is a unique smooth (topological) structure on $S$ such that $f:M\to S$ is a diffeomorphism (homeomorphism).
    And the structure on $S$ is compatible with $N$, meaning that the inclusion $S\hookrightarrow N$ is a smooth (topological) embedding.
\end{remark}

\subsubsection{The Whitney Embedding Theorem}
\begin{theorem}[Whitney Embedding Theorem]
    Every smooth $n$-manifold with or without boundary admits a proper smooth embedding into $\mathbb{R}^{2 n+1}$.
\end{theorem}
\begin{corollary}
    Every smooth n-dimensional manifold with or without boundary is diffeomorphic to a properly embedded submanifold (with or without boundary) of $\mathbb{R}^{2 n+1}$.
\end{corollary}


\begin{corollary}
    Suppose $M$ is a compact smooth $n$-manifold with or without boundary. If $N \geq 2 n+1$, then every smooth map from $M$ to $\mathbb{R}^N$ can be uniformly approximated by embeddings.
\end{corollary}




\begin{theorem}[Whitney Immersion Theorem]
    Every smooth n-manifold with or without boundary admits a smooth immersion into $\mathbb{R}^{2 n}$.
\end{theorem}

\begin{theorem}[Strong Whitney Immersion Theorem]
    If $n>1$, every smooth $n$-manifold admits a smooth immersion into $\mathbb{R}^{2 n-1}$.
\end{theorem}





\section{Submanifold}

\begin{definition}
    An \textbf{embedded submanifold} of $M$ is a subset $S\subseteq M$ such that
    \begin{enumerate}
        \item $S$ is a manifold
        \item the inclusion map $i:S\hookrightarrow M$ is a smooth embedding.
    \end{enumerate}
\end{definition}


\begin{definition}
    Given an subset $S\subset M$, $k\in \mathbb{N}$ and embedding $i:S\hookrightarrow M$.
    \begin{enumerate}
        \item
              If there is a chart $(U,\varphi)$ covering $S$ and $\varphi(S)=\left\{x\in \varphi(U)| x^i=0 ,i=k+1,\ldots,m\right\}$, then we say that $S$ is a \textbf{ $k$-slice} of $U$.
        \item
              we say that $S$ satisfies the \textbf{local $k$-slice condition} if each point of $S$ is contained a chart $(U,\varphi)$ for $M$ such that $S \cap U$ is a $k$-slice
              in $U$.
    \end{enumerate}
    Any such chart is called a \textbf{slice chart} for $S$ in $M$, and the corresponding coordinates $\left(x^1, \ldots, x^n\right)$ are called \textbf{slice coordinates}.
\end{definition}
\begin{theorem}[Local Slice Criterion for Embedded Submanifolds]
    If $S\subset M$ is an embedded $k$-dimensional submanifold, then $S$ satisfies the local $k$-slice condition.
\end{theorem}
\begin{theorem}
    If $M$ is a smooth $n$-manifold with boundary, then $\partial M$ is an embedded $(n-1)$-dimensional submanifold of $M$.
\end{theorem}

\begin{theorem}
    Let $f:M\to N$ be a smooth map with constant rank $r$. Then each level set $f^{-1}(c)$ is a properly embedded submanifold of codimension $r$ in $M$.
\end{theorem}

\begin{corollary}
    Every regular level set of a smooth map is a properly embedded submanifold whose codimension is equal to the dimension of the codomain.
\end{corollary}








\end{document}
